% Round One (NeurIPS poster) submission, Smart Buildings Challenge 2026.
%
% Written against plain article rather than the NeurIPS workshop style file,
% which is not public at the time of writing. Swap \documentclass and the
% geometry line for the official template before submitting; nothing else in
% this file depends on the class.
\documentclass[10pt]{article}

\usepackage[margin=1in]{geometry}
\usepackage{amsmath}
\usepackage{amssymb}
\usepackage{booktabs}
\usepackage{microtype}
\usepackage{url}
\usepackage[colorlinks=true,allcolors=blue]{hyperref}
\usepackage{caption}

\title{\textbf{Airflow Is the Load in Heating and the Delivery in Cooling}\\[2pt]
\large One Asymmetry Decides Three Controller Decisions in the Smart
Buildings Challenge}

\author{%
  Nhat Trinh Ngoc Minh\\[2pt]
  \small Faculty of Information Technology, Van Lang School of Technology\\
  \small Van Lang University, 69/68 Dang Thuy Tram Street,\\
  \small Binh Loi Trung Ward, Ho Chi Minh City, Vietnam\\[2pt]
  \small\texttt{nhat.2474802010280@vanlanguni.vn}\\
  \small ORCID: \href{https://orcid.org/0009-0006-9477-9642}{0009-0006-9477-9642}
}

\date{}

\begin{document}
\maketitle

\begin{abstract}
The Smart Buildings Challenge asks whether a single HVAC control policy can
transfer to buildings and climates it has never seen, and poses as an open
question whether domain-specific models or generic foundation models will
generalize best. We report a result on the domain-specific side, and a caution
about how it should be measured.

One property of the plant decides three separate controller choices, and the
shipped reference controllers get all three wrong in the same direction. The
air handler cannot heat---all heat comes from the VAV reheat coils---and below
about $-6.8^{\circ}$C ambient its supply-air command is inert as well, because
it delivers $\min(\text{mixed air}, \text{setpoint})$ and with recirculation
capped at $70\%$ the mixed air is already below the coldest setpoint
available. So \textbf{in heating, airflow is the load}: three tenths of every
cubic metre the fan moves is ambient air that must then be heated, and the fan
work is spent before the gas is spent undoing it. \textbf{In cooling the same
airflow is the delivery.} Every decision that adds or withholds airflow
therefore pays in one direction and costs in the other, and each needs the same
gate: static pressure, air-handler run length, and whether to run a handler at
all.

Acting on this takes the metric from $0.05496$ to $0.02774$ over the 30
published scenarios, a $49.5\%$ reduction, with no learned parameters and in
pure NumPy. An earlier version of the controller measured $0.03785$ locally
and scored $0.04932$ on the organizers' held-out leaderboard against
$0.06404$ for the baseline, so the gain transfers at $23\%$ and the two scales
are readable against each other.

We also report eight refuted hypotheses, including our own initial one and one
conclusion we drew and then had to withdraw: after six failed attempts at the
comfort term we claimed the shipped gains sat near a local optimum. They do
not. Every one of those attempts pushed the zone loop \emph{harder}, and the
loop was already oscillating; the fix is a five-times softer gain with an
integral term, which no single-knob ablation can find because neither half
works alone.
\end{abstract}

\section{Introduction}

The challenge \cite{goldfeder2026} evaluates one control policy on six unseen
buildings at five unseen weather stations. Buildings differ in floor plan, zone
count, air-handler count and envelope construction, so both the observation and
action vectors change shape per building and there is no fixed layout to
memorize. The organizers state the core hurdle as transferability, and ask
whether ``domain specific models will win out, or whether generic foundation
models will generalize best'' \cite{goldfeder2026}.

Two features of the setup push against learned policies. Per-rollout wall-clock
is bounded at ten times the \texttt{bang\_bang} baseline, which rules out
winning by online planning. And the metric is a mean over 30 equally weighted
scenarios whose difficulty is very unequal: in our measurements the four worst
scenarios contribute $38\%$ of the strongest baseline's total loss. Mean-case
improvement is worth a fraction of tail-case improvement, and tails are where
sample-efficient learning is weakest.

This paper takes the other route. We read the reward function and the plant
model as a price list, find one actuator whose price the reference controllers
have the wrong sign for, and correct it.

\section{The reward as a price list}

Per five-minute step the reward is
\begin{equation}
r_t = 0.6\,p_t - 0.2\,c_t - 0.2\,e_t ,
\end{equation}
with all three terms penalties, $r_t \le 0$, and the leaderboard metric the mean
over scenarios of $|\bar r|$. Substituting the simulator's own cost and carbon
models at the shipped normalisation constants
($320$\,kW electricity, $820$\,kW gas) gives the marginal cost in metric units
of one delivered kilowatt-hour:

\begin{center}
\begin{tabular}{lc}
\toprule
 & metric cost per kWh \\
\midrule
fan / pump electricity            & $0.0064$ \\
gas heat (boiler, $\eta\approx1$) & $0.0034$ \\
chiller cooling (COP $3.5$)       & $0.0018$ per kWh \emph{thermal} \\
\bottomrule
\end{tabular}
\end{center}

Electricity costs $1.9\times$ gas per kWh of fuel, and a fan has no coefficient
of performance to divide by where a chiller has $3.5$. \textbf{Moving air is the
most expensive thing this plant can do.}

The comfort term is an occupancy-weighted mean of per-zone shortfall, where
shortfall is $\mathrm{dev}/5\,\mathrm{K}$ inside the knee and quadratic and
unbounded beyond it. Every occupied zone sitting $1$\,K outside its band
therefore costs $0.6 \times 0.2 = 0.12$ per step, which at $0.0064$ per kWh is
$19$\,kWh per step---$225$\,kW held continuously, against a $320$\,kW plant.
\emph{While comfort is scored, energy is nearly free beside it.}

Finally, the comfort term is gated hard on occupancy:
\verb|total_occupancy > 0.0| or the term is identically zero whatever the zone
temperatures are. Occupants arrive between hours 6 and 12, leave between 13 and
19, and are absent on weekends and holidays, leaving roughly $60\%$ of the
$8{,}640$ scored steps with the comfort term switched off.

\section{Airflow is the load}

The gated comfort term invites an obvious move: stop conditioning an empty
building. We implemented it and it was \emph{worse} than the strongest baseline
on every scenario we tried. Section~\ref{sec:negative} gives the two reasons.
The useful finding came from asking instead why the reference controllers spend
what they spend.

Every policy in the shipped \texttt{pid} family sets static pressure from the
control effort $u$:
\begin{equation}
\texttt{pressure} = \mathrm{clip}(2|u| - 1,\,-1,\,1).
\label{eq:pid}
\end{equation}
The baseline's own documentation defends this against pinning pressure at
maximum, and against that alternative it is an improvement. But $|u|$ does not
distinguish the two directions of the loop, and the directions are not
symmetric.

\paragraph{Heating.} The air handler cannot heat. It mixes return and outside
air, cools the mixture towards the commanded supply temperature, and all heat
in the building comes from the VAV reheat coils fed by the hot-water loop. In
cold weather the supply-air action is inert as well: the handler delivers
$\min(\text{mixed air}, \text{setpoint})$, and with recirculation capped at
$0.70$, return air near $293$\,K and ambient at $231$\,K,
\begin{equation}
\text{mixed} = 0.70 \times 293 + 0.30 \times 231 = 274\ \mathrm{K},
\end{equation}
already below the $285$\,K floor the setpoint can command.

So in heating the fan moves $274$\,K air past a coil, and every cubic metre of
it must be heated from $274$\,K before it does the zone any good. The
ventilation floor makes this unavoidable: three tenths of the moved air is
ambient, at $-42^\circ$C in the coldest shipped strip. \textbf{The fan work is
spent, and then gas is spent undoing the air the fan moved.}

The thermal energy a zone needs is fixed by its heat loss. It can arrive as a
lot of slightly warm air or a little hot air, and the second is strictly cheaper
here, because the gas is the same either way and the fan is not. Equation
\eqref{eq:pid} does the opposite: it buys more of the expensive fuel in order to
import more of the load.

\paragraph{Cooling.} Air is the only delivery mechanism, so flow must be bought.
But the trade still favours temperature over flow---the chiller kilowatt-hour is
a third of the fan kilowatt-hour---and the same ventilation floor means that in
$45^\circ$C ambient, extra flow again imports extra load.

\section{Controller}

Our policy is the shipped \texttt{pid\_water} controller with four changes and
no learned parameters. Three of the four are the same asymmetry applied to a
different actuator.

\paragraph{A flat band.} \texttt{pid\_water} widens its band to an eco setback
outside comfort hours. We hold the comfort band around the clock. This is not
an energy argument but a tail argument: a building held in band overnight starts
the morning with no deviation to recover, and the recovery penalty is larger
than the holding cost. The shipped \texttt{fixed\_setpoint} already hints at
this by beating \texttt{pid\_water} in cold climates.

\paragraph{A direction-aware pressure law.} Where \eqref{eq:pid} applies, we
instead clamp static pressure to its floor whenever the handler is heating
($u>0$), and scale it by $0.6$ when cooling. The simulator enforces a minimum
static pressure under any running handler, so the floor requests that minimum
rather than no air; ventilation is not something a policy may trade away, and
this does not try to.

One further change follows from the same reading. \texttt{pid\_water} opens the
reheat valve in proportion to shortfall, fully open at $2$\,K. That protects a
mild climate from overshoot and starves a cold one, where the valve is the only
heat source and the air reaching it is already $20$\,K below the zone. We
saturate the valve at $0.5$\,K.

\paragraph{A minimum run length for cooling runs.} The shipped loop idles a
handler the instant its \emph{aggregate} zone error clears and restarts it at
the next excursion. Because the zone loop oscillates (below), those excursions
arrive every few steps: every handler switches on or off 50 to 107 times a
day. Each restart is paid for twice---the fan spins up against a cold duct, and
the zones drift while it is off, so the restart is larger than the excursion
that triggered it. We hold a run that was started to cool at its last commands
for 288 steps.

Heating runs are not held, and that is the asymmetry again: a handler that
cannot heat has no reason to stay on, and holding one in the cold is a
continuous flow of $274$\,K air for nothing. We also tested forcing every
handler on through the comfort window---what a real VAV plant does, supplying
cool air steadily while each box decides how much to take. Gated on the zones'
own demand it is worth $-44\%$ on \texttt{phoenix\_4} and exactly $0.0\%$ on
\texttt{fairbanks\_1}, because the gate never fires there. Ungated it costs
$+54\%$ in Fairbanks. But combined with the hold it is worse than the hold
alone ($0.0506$ against $0.0385$): both mechanisms exist to stop a handler
being cycled off, so together they interfere. We kept the hold.

\paragraph{Per-zone PI against the band centre.} The shipped zone law is
\begin{equation}
\texttt{damper} = \min\!\left(1, \frac{|\mathrm{deadband\_error}|}{2\,\mathrm{K}}\right),
\end{equation}
proportional only, and aimed at the band \emph{edge}: \texttt{deadband\_error}
is zero everywhere inside the band, so the loop has no target in there to
hold. It pushes until the zone crosses the edge, stops, and the zone drifts
back out the other side. Zones average \emph{inside} the $21$--$24^{\circ}$C
band and yet sit outside it $57$--$61\%$ of working hours, swinging $2$--$4$\,K
every few steps \cite{probezt}.

We give each zone its own PI loop on $e = T - \text{centre}$, with
$k_p=0.1$, $k_i=0.005$ per five-minute step and conditional integration, from
an hour before the comfort window until it closes. One signal drives both
actuators: the damper follows $|u|$, because air is needed either to cool or to
carry the reheat coil's output, and only the valve follows the sign.

\section{Results}

All runs are full $8{,}640$-step episodes. The simulator is deterministic---we
verified two runs of one policy on one scenario reproduce to ten decimal
places---so each number is a single rollout and no averaging is required.

\begin{table}[h]
\centering
\caption{Leaderboard metric (lower is better) over the 30 published training
scenarios, with the weighted contribution of each reward term. Both policies
re-run with the current simulator.}
\begin{tabular}{lccccc}
\toprule
policy & metric & worst scenario & comfort & energy & carbon \\
\midrule
\texttt{pid\_water} (strongest baseline) & $0.05496$ & $0.21392$ & $0.03489$ & $0.01150$ & $0.00857$ \\
band + fan clamp only & $0.03785$ & $0.07262$ & $0.02655$ & $0.00678$ & $0.00452$ \\
\textbf{+ zone PI + cooling hold} & $\mathbf{0.02774}$ & $\mathbf{0.07685}$ & $0.01576$ & $0.00706$ & $0.00492$ \\
\midrule
change & $-49.5\%$ & $-64.1\%$ & $-54.8\%$ & $-38.6\%$ & $-42.6\%$ \\
\bottomrule
\end{tabular}
\label{tab:main}
\end{table}

\textbf{The fan clamp improves all three terms at once}, which is the
signature of a mispriced actuator rather than a trade: less static pressure is
less airflow, which is less ambient air imported, which is less load, so the
coils keep up where before they could not. Comfort improves \emph{because} the
energy was cut, not despite it.

The zone PI and the hold then take comfort down by a further $41\%$ for a
slight rise in energy, which is the trade the $35{:}1$ price ratio says to
make.

\textbf{The remaining tail is now entirely one climate.} The four worst
published scenarios are \texttt{fairbanks\_1}, \texttt{fairbanks\_7},
\texttt{fairbanks\_8} and \texttt{fairbanks\_2}, at $0.077$, $0.067$, $0.061$
and $0.060$. The held-out \texttt{dev} split says the same thing across
independent buildings, by station:

\begin{center}
\begin{tabular}{lccccc}
\toprule
 & atlanta & phoenix & singapore & denver & fairbanks \\
\midrule
\texttt{dev} score & $0.0137$ & $0.0137$ & $0.0152$ & $\mathbf{0.0655}$ & $\mathbf{0.0644}$ \\
\bottomrule
\end{tabular}
\end{center}

The two cold stations are four to five times the other three. That is the
design showing through rather than a surprise: both of the later changes are
gated on cooling, so neither fires in the cold, and what remains is exactly
the climate they do not touch. It is also where the asymmetry says the problem
is hardest---in heating the one actuator that moves air is working against the
load it is trying to serve, so there is no cheap direction to push.

The gain is concentrated in the tail, which is what the metric rewards. Our
policy is better on 19 of 30 scenarios; summed over scenarios it improves
$0.575$ and regresses $0.031$, a ratio of $19{:}1$, and the median regression
among the eleven is $0.001$. The four largest improvements are
\texttt{fairbanks\_1} $0.214\!\to\!0.072$, \texttt{phoenix\_4}
$0.160\!\to\!0.043$, \texttt{phoenix\_5} $0.136\!\to\!0.049$ and
\texttt{singapore\_4} $0.110\!\to\!0.024$: the two coldest and the two hottest,
which is what a mechanism acting through the ventilation floor predicts.

\begin{table}[h]
\centering
\caption{Ablation, mean over \texttt{fairbanks\_1}, \texttt{singapore\_6} and
\texttt{phoenix\_4}. Each row adds one change to the row above.}
\begin{tabular}{lcl}
\toprule
 & metric & \\
\midrule
\texttt{pid\_water}                     & $0.13478$ & shipped baseline \\
\quad + flat band                       & $0.12321$ & band alone \\
\quad + cooling clamp only              & $0.11315$ & \\
\quad + heating clamp only              & $0.05226$ & \textbf{the single biggest term} \\
\quad + reheat saturation at $0.5$\,K   & $0.05034$ & \\
\quad + per-zone PI                     & $0.04657$ & \\
\quad + cooling-run hold, 288 steps     & $\mathbf{0.03846}$ & final \\
\midrule
\multicolumn{3}{l}{\emph{tested and dropped}} \\
\quad PI, proportional only ($k_i=0$)   & $0.05304$ & worse than no PI at all \\
\quad PI + forced-on handlers           & $0.04248$ & \\
\quad PI + forced-on + hold             & $0.05061$ & the two interfere \\
\bottomrule
\end{tabular}
\label{tab:ablation}
\end{table}

\subsection{Held-out buildings and climates}

Because the published 30 scenarios are a training set, we built two
replacements from material they never use. The generator is procedural, ten of
the 48 shipped weather strips are unused by the published set, and a further
station, \texttt{minneapolis}, is used by none of them. \textbf{dev} is four
newly generated buildings crossed with five unused strips at seen stations;
\textbf{test} is four further buildings crossed with \texttt{minneapolis}, so
neither the building nor the climate has been seen. All eight fall inside the
$28$--$85$ zone and $3$--$9$ air-handler range the organizers state for the
evaluation set.

\begin{table}[h]
\centering
\caption{The same controller on three sets. The gain is positive everywhere
and much smaller off the published scenarios.}
\begin{tabular}{lcccc}
\toprule
 & \texttt{pid\_water} & fan clamp & full & change \\
\midrule
published (30 scenarios)         & $0.05496$ & $0.03785$ & $0.02774$ & $-49.5\%$ \\
dev (20, unseen buildings)       & $0.04630$ & $0.04336$ & $0.03449$ & $-25.5\%$ \\
test (8, unseen also climate)    & $0.11253$ & $0.09781$ & $0.08937$ & $-20.6\%$ \\
\midrule
\textbf{held-out leaderboard}    & $\mathbf{0.06404}$ & $\mathbf{0.04932}$ & --- & $\mathbf{-23.0\%}$ \\
\bottomrule
\end{tabular}
\label{tab:holdout}
\end{table}

The last row is the organizers' own held-out set, obtained by submitting the
unmodified \texttt{pid\_water} baseline alongside our entry. A baseline whose
local value is known exactly turns the leaderboard into a ruler: it scores
$1.17\times$ \emph{higher} there than on the published scenarios, so the
evaluation set is slightly harder rather than easier, and the $23\%$ gain is
measured on buildings and weather we never saw.

Our \texttt{test} split is the weakest of the three estimates and we would not
lean on it: it is entirely \texttt{minneapolis}, a cold station, and both of
the later changes are gated on cooling. It measures the one climate they do
not address, which makes $-20.6\%$ a lower bound rather than a representative
figure.

\textbf{We think the gap is the finding, not a flaw in it.} The worst-case
column explains it: \texttt{pid\_water}'s worst published scenario scores
$0.214$, its worst dev scenario $0.089$. The published 30 contain catastrophic
baseline failures---\texttt{fairbanks\_1}, \texttt{phoenix\_4},
\texttt{phoenix\_5}, \texttt{singapore\_4}---and our change is largely
repairing those. Where the baseline does not fail badly there is less to win.
A method evaluated only on the published set would therefore report roughly
five times the improvement it should expect on fresh buildings from the same
generator.

The decomposition shifts too. On dev, energy and carbon still fall sharply
($0.0073\!\to\!0.0053$ and $0.0071\!\to\!0.0052$) but comfort rises slightly
($0.0318\!\to\!0.0332$), where on the published set all three improved.
\texttt{phoenix} is the only station that regresses, at $+2.9\%$, which points
at the cooling clamp rather than the heating clamp as the term trading comfort
away---consistent with the mechanism, since in cooling the air is the delivery
path and not only the load.

We note one way our holdout may differ from the organizers'. The shipped floor
generator's default footprint is $400$--$700 \times 300$--$500$\,ft, which
produces buildings far outside the stated evaluation range---one seed gave 162
zones and 17 air handlers. We therefore drew footprints from the shipped six's
own range and rescaled until the realised zone count landed in $28$--$85$. If
the organizers' held-out buildings were drawn differently, our estimate of the
transfer gap moves with them.

\section{What did not work}
\label{sec:negative}

We report these because each was a reasonable inference from the reward and each
was wrong, which is itself information about the benchmark.

\paragraph{Switching the plant off when the building is empty.} Our original
hypothesis, and the sign was backwards. Two independent reasons. First, being
empty \emph{inside} the scored window is not an opportunity: arrivals are
independent stochastic draws, so the building is routinely empty for a few steps
after hour 6, and because regret is normalized by \emph{total} occupancy the
first person through the door carries their zone's entire comfort score alone.
Second, a controller that coasts must later recover, and the recovery lands in
the low-occupancy morning where each occupant's weight is largest. Driving to
the comfort band overnight to avoid that cost \emph{more} than coasting saved:
unoccupied-step loss rose from $447$ to $830$ for a comfort gain of $0.010$.

\paragraph{Deeper setbacks.} Monotonically worse: $0$\,K $0.123$, $3$\,K
$0.166$, $5$\,K $0.160$, $8$\,K $0.183$. Steady holding beats cycling here.

\paragraph{Aiming inside the comfort band.} The penalty is flat zero anywhere
inside the band, so a margin looked like free insurance against drift. It is
clearly harmful ($0.236$ against $0.123$), and harmful on the \emph{comfort}
term, not merely the energy term. Narrowing a $3$\,K band to $2$\,K leaves the
loop, which saturates at $1/k_p = 2$\,K, oscillating across the real edge.

\paragraph{Forcing the water loop hotter.} Having removed the airflow, raising
the water temperature to compensate looked necessary. It is not: carbon rose
from $0.0134$ to $0.0162$. With the flow down, the baseline's own demand reset
is already hot enough.

\paragraph{Suspending the fan economy when comfort is at stake.} Against the
$35{:}1$ price ratio this should pay---a quiet fan delivers heat slowly, and
slow is only cheap when nobody is being charged. It does not pay ($0.074$
against $0.050$), and the reason is the mechanism itself: the fan is what
\emph{causes} the load, so spending the expensive fuel to deliver heat imports
more cold air than the extra flow delivers.

\paragraph{Capping the hot-water pump.} Also electricity, also pinned at
maximum by the baseline, and also worse when capped ($0.051$ against $0.050$).

\paragraph{Everything aimed at the comfort term, and a conclusion we withdrew.}
After the fan clamp, comfort is $70$--$75\%$ of what is left on all three sets.
Against a reference of $0.0503$ on the three hard scenarios we measured:
reading each zone's band from its own setpoint sensors $0.0548$; a stiffer
loop, $k_p$ $0.5\!\to\!1.0$, $0.0794$; dampers opening at $1$\,K instead of
$2$\,K $0.0897$, and at $0.5$\,K $0.1117$; driving the air-handler loop from
its worst zone rather than the mean $0.0970$; a per-zone damper gate
withholding airflow from a zone the supply air cannot help $0.0626$; and a
margin aiming inside the band rather than at its edge, $0.236$.

Six attempts, all worse, and we concluded that the shipped gains sat near a
local optimum and the comfort term was not reachable by retuning.
\textbf{That conclusion was wrong}, and it is worth reporting why, because the
error is structural rather than careless.

Four of the six---margin, $k_p$, damper gain, worst-zone---are pushes in the
direction that \emph{causes} the oscillation, applied to a loop that was
already oscillating. The correct direction was five times softer. But a softer
proportional gain alone is also worse: measured, $k_p=0.1$ with $k_i=0$ scores
$0.0530$ against $0.0503$ for leaving the loop alone, because soft P trades
the swing for a standing offset. The gain and the integral only work together,
and aiming at the band centre rather than the edge is required for the
integral to have a target at all.

\textbf{A one-factor-at-a-time ablation cannot find a change whose two halves
are each harmful alone.} That is the methodological lesson, and six consistent
negative results made it look like evidence about the problem when it was
evidence about the search.

Two of the six do answer questions the starter kit leaves open: it describes
both the sensor band and the worst-zone loop in its documentation and ships
with both switched off. Our measurements say it was right to.

\section{A note on the shipped baseline table}

The starter kit ships \texttt{baseline\_returns.json}, and
\texttt{score\_leaderboard.py} prints it beside a participant's score so the
four baselines need not be re-run. We find it no longer reproduces:
\texttt{pid\_water} on \texttt{building\_train\_1/fairbanks\_1} gives
$0.21392$ with current code against the recorded $0.21168$, and $0.05496$
against $0.05414$ over the 30 scenarios.

The file was committed with the initial starter kit and never regenerated,
while a later commit fixed the simulation tracker querying occupancy with VAV
ids---a fix whose own message notes that it ``creates fictitious populations
and consumes random draws, changing subsequent simulation trajectories''. The
trajectories moved and the table did not. The discrepancy is $1.1$--$1.5\%$,
which is small but is the same order as several of the tuning effects above.
Participants comparing against the shipped table are measuring against a
simulator that no longer exists.

\section{Discussion}

On the organizers' open question, this is evidence that domain structure is
worth a good deal here, and cheaply. The result is a 40-line change,
dimensionless, with no fitted constants and nothing to transfer, and its sign
is the same on every set we measured because it acts through a
constraint---the $70\%$ recirculation cap---that is a property of the plant
rather than of any building.

The headroom here divides into two kinds, and they are found in different
ways. One is \emph{pricing}: the controller buying the most expensive fuel in
order to import the load it was trying to remove. That one is visible by
reading the plant and the reward, and it is what the heating-side fan clamp
corrects. The other is \emph{tracking}: a zone loop oscillating across the band
it is meant to hold, for $57$--$61\%$ of working hours. That one is visible
only by logging what the controller actually does, step by step, which is how
it was found \cite{probezt} and not how we were looking.

We spent six experiments concluding the second kind did not exist, on the
strength of six failures that were all the same mistake. The honest summary is
that reading the reward told us where the money was, and it was not sufficient
to tell us how to collect it.

Its \emph{magnitude}, though, is a property of the evaluation set, and that is
the part we would most want other entrants to take from this. The difference
between $33\%$ and $5.7\%$ is not noise and not overfitting in the usual
sense---nothing was tuned on the published scenarios that is not also present
on dev---but a difference in how often the baseline's failure mode occurs. Any
entry that reports only published-set numbers is reporting the frequency of
those failures as much as the quality of its method.

We would not read it as evidence against learned policies in general. We would
read it as a statement about what a learned policy would have to discover: not a
smooth function of the observation, but the sign of a coupling between an
actuator and the load it is meant to serve, visible only in the regime where the
supply-air action has gone inert. A policy-gradient method optimizing an
$8{,}640$-step return at ten times \texttt{bang\_bang}'s runtime budget has a
hard road to that. The natural synthesis is to let structure carry the
mechanism and let learning carry what structure cannot express---which is the
direction we intend for the final submission.

\paragraph{Limitations.} The tuning in Table~\ref{tab:ablation} beyond the
heating clamp is small, measured on three scenarios, and we do not claim those
constants transfer. The cooling clamp in particular is a wash once measured
properly: dropping it moves the score by $+2.7\%$ on the published set,
$-0.7\%$ on dev and $+0.3\%$ on test, which averages to $0.6\%$ and is not a
difference we would defend. \textbf{The heating clamp is the result; the rest
is noise around it.} Our policy is better on 19 of 30 published scenarios and only 8 of
20 dev scenarios; it wins on the mean because the losses are tiny (median
$0.001$) and the wins are large, which is the right shape for this metric but
is not the same as being uniformly better. Our held-out buildings are our own
construction and may not match the organizers' draw. We have not run against
the organizers' held-out set, and the public leaderboard scores only a
rotating subset of it.

\paragraph{Reproducibility.} Code, every sweep in this paper as JSON, and the
held-out generator are at
\url{https://github.com/MinhNhat-2504/hvac-generalization-lab}. Each negative
result above is a registered policy variant, runnable by name.

\begin{thebibliography}{9}
\bibitem{goldfeder2026}
J.~Goldfeder, A.~Rai, P.~M.~Wyder, N.~Li, V.~Dean, D.~Vakalis, J.~Lin, B.~Dong,
M.~Rossetti, N.~Kutz, S.~Arisaka, J.~Sipple, H.~Lipson, and Y.~Bengio.
\newblock The Smart Buildings Challenge: Learning Foundation Models for
Real-World Optimization at Scale.
\newblock arXiv:2610.18511, 2026.

\bibitem{starterkit}
Smart Buildings 2026 Organizers.
\newblock Smart Buildings Competition Starter Kit.
\newblock \url{https://github.com/smart-buildings-2026/starter-kit}, 2026.

\bibitem{probezt}
Smart Buildings: per-zone PI control and AHU cooling hold.
\newblock Public competition notebook, Kaggle Code tab, 2026.
\newblock The per-zone oscillation diagnosis, the PI gains and the cooling-run
hold are from this notebook; our implementations are independent. It also
reports the minimum-fan-pressure result we arrived at separately, with the
same explanation, which we take as mutual confirmation of the mechanism.
\end{thebibliography}

\end{document}
